\section[Cote 158 : la proposition 1 sur les monoïdes finis]{Cote 158 : la proposition 1 sur les monoïdes finis\protect\verifie{Folder158}}\label{sec:158}

La cote 158, \emph{[Autour des Dérivateurs] : notes manuscrites (s.d.)}, quatre-vingt-cinq pages que l'inventaire date \q{[vers 1990-1991]}, étudie une seule question : une petite catégorie $X$ est-elle cofiltrante dès que tout objet $0$-connexe du topos des préfaisceaux sur $X$ est $1$-connexe ? Les pages~44 à~53, sous sa pagination~1 à~8, traitent le cas fini, en commençant par celui d'un monoïde, c'est-à-dire d'une catégorie à un objet. La proposition~1 est énoncée page~44 et démontrée jusqu'au haut de la page~46, qui conclut \q{cqfd} (\lot{158}{3}{44}). Sous l'énoncé, un premier paragraphe, barré de deux longues diagonales, commence par \q{L'unicité de $p$ est claire}. La démonstration reprend ensuite par la pseudo-cofiltrance et réalise $M$ comme monoïde d'endomorphismes d'un ensemble $E$, le mot \q{fini} devant $E$ étant biffé. Elle écrit : \q{comme $E_0$ est fini, on en conclut que $u(E_0) = E_0$} (\lot{158}{3}{45}).

\begin{definition}
Un monoïde $M$ est \emph{pseudo-cofiltrant} si, pour tous $u, v \in M$, il existe $u', v' \in M$ tels que $uu' = vv'$ ; il est \emph{cofiltrant} s'il est pseudo-cofiltrant et si, pour tous $u, v \in M$, il existe $w \in M$ tel que $uw = vw$. Le \emph{groupe enveloppant} de $M$ est \emph{trivial} si tout morphisme de monoïdes de $M$ dans un groupe est trivial. La première définition est celle que la lecture tire de la page~55, où le manuscrit en donne l'équivalent ; la dernière est la propriété universelle du groupe enveloppant.
\end{definition}

\begin{enonce}[Proposition 1, pages 44 à 46\piste{158-finite-monoid-right-zero}]
Soit $M$ un monoïde fini dont le groupe enveloppant est trivial et qui est pseudo-cofiltrant. Alors il existe $p \in M$ tel que $up = p$ pour tout $u \in M$ --- en particulier $p^2 = p$ --- et $M$ est cofiltrant.
\end{enonce}

\begin{lemme}\label{lem:158-etape1}
Dans un monoïde fini pseudo-cofiltrant $M$, il existe $p$ tel que, pour tout $u \in M$, on ait $p = uv$ pour un $v \in M$ : $p$ appartient à tous les idéaux à droite principaux $uM$.
\end{lemme}

\begin{proof}
Montrons, par récurrence sur le cardinal d'une partie finie $F \subseteq M$, qu'il existe $p$ appartenant à $uM$ pour tout $u \in F$ ; appliqué à $F = M$, c'est l'énoncé. Pour $F$ vide, $p = 1$ convient. Supposons l'assertion vraie pour $F$, avec un élément $p$, et soit $a \in M$. Par pseudo-cofiltrance appliquée à $p$ et $a$, il existe $p', v'$ tels que $pp' = av'$. L'élément $pp'$ est dans $aM$ ; et pour $u \in F$, si $p = uv$, alors $pp' = u(vp')$ est dans $uM$. L'élément $pp'$ convient donc pour $F \cup \{a\}$.
\end{proof}

\begin{proof}[Démonstration de la proposition 1]
Soit $p$ donné par le lemme~\ref{lem:158-etape1}, et $E_0 = pM$ l'ensemble des $py$, $y \in M$. C'est l'image de $p$ dans la représentation régulière gauche de $M$, où $u$ opère sur $M$ par $x \mapsto ux$, et $E_0$ est fini puisque $M$ l'est.

\emph{Chaque $u$ envoie $E_0$ sur lui-même.} Soit $u \in M$. D'abord $E_0 \subseteq uE_0$ : le lemme appliqué à l'élément $up$ fournit $v$ tel que $p = upv$, et pour $y \in M$, $py = upvy = u\bigl(p(vy)\bigr)$ est dans $uE_0$. Ensuite $uE_0$, image de $E_0$ par une application, a au plus autant d'éléments que $E_0$ ; une partie de l'ensemble fini $uE_0$ qui a au moins autant d'éléments que lui lui est égale, et donc $uE_0 = E_0$. L'application $x \mapsto ux$ est ainsi une surjection de l'ensemble fini $E_0$ sur lui-même, donc une bijection.

\emph{Le morphisme vers un groupe.} Notons $\varphi(u)$ la permutation $x \mapsto ux$ de $E_0$. On a $\varphi(1) = \mathrm{id}$ et, par associativité, $\varphi(uu') = \varphi(u) \circ \varphi(u')$ : $\varphi$ est un morphisme de monoïdes de $M$ dans le groupe des permutations de $E_0$. Le groupe enveloppant de $M$ étant trivial, $\varphi(u) = \mathrm{id}$ pour tout $u$. En l'appliquant à l'élément $p = p \cdot 1$ de $E_0$, on obtient $up = p$ pour tout $u \in M$.

\emph{Conclusion.} Pour $u = p$, $p^2 = p$. Enfin $M$ est pseudo-cofiltrant par hypothèse, et pour $u, v \in M$, $w = p$ vérifie $uw = p = vw$ : $M$ est cofiltrant.
\end{proof}

\begin{remarque}
La finitude sert deux fois : au lemme~\ref{lem:158-etape1}, pour passer de deux idéaux $uM$ à tous, et pour changer l'inclusion $E_0 \subseteq uE_0$ en égalité. Le \q{comme $E_0$ est fini} de la page vaut parce que $E_0 = pM$ est contenu dans $M$. Il ne vaudrait pas pour un ensemble $E$ quelconque, devant lequel la page a biffé \q{fini}.
\end{remarque}

\begin{remarque}
Sans la pseudo-cofiltrance, l'énoncé est faux : le monoïde $\{1, a, b\}$ où $ax = a$ et $bx = b$ pour tout $x$ est fini, de groupe enveloppant trivial, et n'a pas d'élément $p$ tel que $ap = p = bp$. Sans la finitude, il est faux pour le monoïde $\mathrm{Ep}(E)$ des surjections d'un ensemble infini $E$, que la lecture étudie à sa section~6. (Non démontré ici.)
\end{remarque}

\begin{remarque}
L'élément $p$ n'est pas unique, et le noyau de Rees $MpM = pM$ n'est pas nécessairement réduit à un point. Soit $M = \{1, a, b\}$ avec $x \cdot 1 = x$, $x \cdot a = a$, $x \cdot b = b$ pour tout $x$ ; cette loi est associative, d'unité~$1$, pseudo-cofiltrante ($ua = va$ pour tous $u, v$), de groupe enveloppant trivial ($g(a) = g(ba) = g(b)g(a)$ force $g(b) = 1$, et de même $g(a) = 1$) ; on a $ua = a$ et $ub = b$ pour tout $u$, aucun des deux n'est un zéro bilatère ($ab = b$), et l'intersection des idéaux $MxM$ est $\{a, b\}$.
\end{remarque}

\begin{remarque}
La relation $up = p$ pour tout $u$ fait de $p$ un zéro à droite, et l'idéal à gauche $Mp = \{p\}$ est minimal. Une glose qui dirait $p$ \q{zéro à gauche} et le noyau de Rees \q{ponctuel} serait fausse. Le paragraphe barré de la page~44 affirmait l'unicité de $p$, en identifiant $pp'$ à $p'p$ ; l'exemple précédent la réfute.
\end{remarque}

\begin{remarque}
La proposition~1, avec sa réciproque, est proposée à mathlib : pour un monoïde fini, l'existence d'un zéro à droite équivaut aux deux conditions de la page (leanprover-community/mathlib4\#44702).
\end{remarque}
